\section{Ablation studies}
\label{sec:ablations}

\subsection{Number of MC Dropout passes}

We vary the number of stochastic forward passes $T \in \{5, 10, 20, 30\}$
and measure the effect on both segmentation quality and uncertainty
discriminativeness.

\begin{table}[t]
\centering
\caption{Effect of MC Dropout pass count $T$ on DRIVE test set. Dice
and AUC are nearly constant across $T$, confirming that the mean prediction
stabilizes quickly. Unc-AUROC increases from $T = 5$ to $T = 10$ and
plateaus after $T = 20$. Runtime scales linearly with $T$.}
\label{tab:mc_passes}
\small
\setlength{\tabcolsep}{4.6pt}
\resizebox{\textwidth}{!}{%
\begin{tabular}{ccccccc}
\toprule
$T$ & \textbf{Dice} & \textbf{AUC} & \textbf{ECE} & \textbf{Unc-AUROC} &
\textbf{Time (s)} & $\Delta$\textbf{Unc-AUROC vs.\ $T$=30} \\
\midrule
5  & 0.790 & 0.959 & 0.045 & 0.811 & 0.24 & $-$0.027 \\
10 & 0.792 & 0.960 & 0.042 & 0.836 & 0.47 & $-$0.002 \\
20 & 0.791 & 0.960 & 0.043 & 0.832 & 0.93 & $-$0.006 \\
30 & 0.791 & 0.960 & 0.044 & 0.838 & 1.42 & --- \\
\bottomrule
\end{tabular}%
}
\end{table}

The mean prediction (and therefore Dice/AUC) stabilizes by $T = 5$: all
values are within 0.002 Dice. Segmentation quality does not benefit from
more passes.

Uncertainty quality does benefit, but with diminishing returns. The jump
from $T = 5$ to $T = 10$ adds 2.5 percentage points of Unc-AUROC (0.811
to 0.836). From $T = 10$ to $T = 30$, the gain is only 0.2 percentage
points. Given that runtime scales linearly (0.24\,s to 1.42\,s), $T = 10$
offers the best trade-off for latency-sensitive applications.

$T = 30$ remains our default because the absolute cost (1.4\,s per image)
is acceptable and the marginal uncertainty quality gain, while small, is
real. For real-time applications, $T = 10$ loses negligible uncertainty
quality at one-third the cost.

\subsection{Ensemble size}

We train ensembles of $N \in \{2, 5\}$ models.

\begin{table}[t]
\centering
\caption{Effect of ensemble size on DRIVE test set. Moving from 2 to 5
members improves Unc-AUROC by 3 percentage points but nearly triples
inference time.}
\label{tab:ensemble_size}
\small
\setlength{\tabcolsep}{4.6pt}
\resizebox{\textwidth}{!}{%
\begin{tabular}{cccccc}
\toprule
$N$ & \textbf{Dice} & \textbf{AUC} & \textbf{ECE} &
\textbf{Unc-AUROC} & \textbf{Time (s)} \\
\midrule
2 & 0.768 & 0.946 & 0.032 & 0.803 & 1.37 \\
5 & 0.771 & 0.947 & 0.031 & 0.833 & 3.90 \\
\bottomrule
\end{tabular}%
}
\end{table}

Five-member ensembles improve Unc-AUROC by 3 percentage points over
2-member ensembles (0.833 vs.\ 0.803), with minimal Dice improvement.
The cost is $2.9\times$ inference time (3.90\,s vs.\ 1.37\,s) and
$2.5\times$ training cost.

Comparing ensembles to the other methods: a 5-member ensemble (Unc-AUROC
0.833, 3.90\,s) is outperformed by TTA (Unc-AUROC 0.881, 0.13\,s) by a
wide margin in both uncertainty quality and speed. Ensembles do achieve
the best ECE (0.031) among all methods, confirming earlier findings that
ensembles are well calibrated \citep{lakshminarayanan2017simple}, but
this calibration advantage does not translate to better deferral performance.

\subsection{Deferral threshold sensitivity}

We evaluate how sensitive each deferral strategy is to its threshold
parameter by measuring error reduction across a range of deferral rates.

\begin{table}[t]
\centering
\caption{Error reduction at fixed deferral rates (5\%, 10\%, 20\%, 30\%)
for each combination of uncertainty method and deferral strategy. Values
are pixel-wise error reduction ratios. Confidence-aware deferral
consistently dominates at low deferral rates.}
\label{tab:threshold_sensitivity}
\small
\setlength{\tabcolsep}{4.6pt}
\resizebox{\textwidth}{!}{%
\begin{tabular}{llcccc}
\toprule
\textbf{Uncertainty} & \textbf{Deferral} &
\textbf{5\%} & \textbf{10\%} & \textbf{20\%} & \textbf{30\%} \\
\midrule
MC Dropout & Global     & 8.2\%  & 14.7\% & 23.5\% & 28.1\% \\
MC Dropout & Conf-aware & 21.4\% & 38.6\% & 52.1\% & 57.3\% \\
\midrule
TTA        & Global     & 15.3\% & 28.1\% & 44.8\% & 56.2\% \\
TTA        & Adaptive   & 18.7\% & 35.9\% & 61.3\% & 75.4\% \\
TTA        & Conf-aware & 24.8\% & 43.2\% & 60.7\% & 68.1\% \\
\bottomrule
\end{tabular}%
}
\end{table}

At a 10\% deferral rate (a realistic operating point for clinical workflows),
TTA with confidence-aware deferral reduces errors by 43.2\%, compared to
14.7\% for MC Dropout with global thresholding. The 3$\times$ difference
in error reduction at the same deferral budget is the most practically
relevant comparison in this paper.

Adaptive thresholding dominates at high deferral rates (20--30\%) because
its per-image normalization ensures that even easy images contribute to
deferral. Confidence-aware deferral dominates at low deferral rates (5--10\%)
because it concentrates deferrals on the highest-risk pixels. In a
workflow where review capacity is limited, confidence-aware deferral
is the better choice.

\subsection{Cross-validation stability}

We run 5-fold cross-validation on the 20 DRIVE training images to assess
the stability of our metrics.

\begin{table}[t]
\centering
\caption{Five-fold cross-validation results on DRIVE (mean $\pm$ std).
Low variance across folds indicates stable training and evaluation.}
\label{tab:cv}
\small
\setlength{\tabcolsep}{4.6pt}
\resizebox{\textwidth}{!}{%
\begin{tabular}{lc}
\toprule
\textbf{Metric} & \textbf{Mean $\pm$ Std} \\
\midrule
Dice         & 0.780 $\pm$ 0.004 \\
AUC          & 0.963 $\pm$ 0.003 \\
ECE          & 0.044 $\pm$ 0.006 \\
Unc-AUROC    & 0.768 $\pm$ 0.057 \\
Runtime (s)  & 1.52 $\pm$ 0.15 \\
\bottomrule
\end{tabular}%
}
\end{table}

Dice and AUC are stable (coefficient of variation $<$ 1\%). ECE has moderate
variance (CV 14\%), reflecting sensitivity to the specific train/val split,
which is expected given the small dataset. Unc-AUROC has the highest
variance (CV 7.4\%), ranging from 0.71 to 0.83 across folds. This
variability comes from the interaction between the learned decision
boundary and the uncertainty threshold: small changes in which images
are in the validation set can shift the optimal threshold noticeably.

Despite this fold-level variance, the qualitative conclusions (TTA $>$ MC
Dropout, confidence-aware $>$ global) hold across all folds.
